% 2010 年广东卷文科第 9 题选项图：四个正视图（A）—（D）。
% 来源 PDF：普通高考/2010/2010广东文.pdf，SHA-256 dd04b2053a22c854a9abe1ee03241fc04d5a25c91cda1d7997183bdfaa4a8c6f。
% 原页：第 1 页（A3 三栏）第 2 栏上部，书页码 595。
% 重绘依据：D 为正确视图：左端低、右端较高、中央最高且有中央虚线；
%   A、B、C 为干扰项（A 左高右低，B 单调上升，C 单调下降）。
\begin{tikzpicture}[semithick, scale=0.85, every node/.style={inner sep=1.2pt, font=\small}]
  % (A)
  \begin{scope}
    \draw (0,0) -- (3,0) -- (3,0.9) -- (1.5,2.1) -- (0,1.2) -- cycle;
    \draw[dashed] (1.5,0) -- (1.5,2.1);
    \node at (1.5,-0.55) {（A）};
  \end{scope}
  % (B)
  \begin{scope}[shift={(5,0)}]
    \draw (0,0) -- (3,0) -- (3,2.1) -- (0,0.9) -- cycle;
    \draw (0,0.9) -- (3,1.2);
    \node at (1.5,-0.55) {（B）};
  \end{scope}
  % (C)
  \begin{scope}[shift={(10,0)}]
    \draw (0,0) -- (3,0) -- (3,1.2) -- (0,2.1) -- cycle;
    \draw (0,1.2) -- (3,0.9);
    \node at (1.5,-0.55) {（C）};
  \end{scope}
  % (D)
  \begin{scope}[shift={(15,0)}]
    \draw (0,0) -- (3,0) -- (3,1.5) -- (1.5,2.1) -- (0,0.6) -- cycle;
    \draw[dashed] (1.5,0) -- (1.5,2.1);
    \node at (1.5,-0.55) {（D）};
  \end{scope}
\end{tikzpicture}
